\chapter{CPU architecture}
\label{ch:cpuarch}

\section{Execution model}

The core executes RV32I in three stages. The stage split is the one the brief
fixed, but the placement of work inside it is a design decision with a
consequence that runs through the whole implementation:

\begin{table}[H]
\centering
\caption{Pipeline stages and the work done in each}
\label{tab:stages}
\begin{tabular}{@{}clL{0.56\linewidth}@{}}
\toprule
\textbf{Stage} & \textbf{Name} & \textbf{Work} \\
\midrule
S1 & Fetch & PC generation, \sig{ibus} AXI read, branch prediction lookup. \\
S2 & Decode + Execute & Decode, register read, ALU, branch resolution, CSR access,
     the entire load/store transaction, and every trap and interrupt decision. \\
S3 & Writeback & Register file write. \\
\bottomrule
\end{tabular}
\end{table}

Putting decode, execute and the memory access in one stage means an instruction
either completes entirely or not at all. That is what makes the traps precise
without any recovery mechanism: there is never a half-finished instruction for a
trap to repair. The cost is a long combinational path (decode $\to$ ALU $\to$
address $\to$ alignment check), which would limit clock frequency on real
silicon. For the behavioural simulation target of this project that cost is not
paid, so the trade was taken deliberately.

\begin{figure}[H]
\centering
\resizebox{\linewidth}{!}{%
\begin{tikzpicture}[x=1mm,y=1mm,font=\small]
  % --- stage bands -------------------------------------------------------
  \begin{scope}[on background layer]
    \fill[boxbg]     (-16,-64) rectangle (22,12);
    \fill[boxbg!55]  (26,-64)  rectangle (128,12);
    \fill[boxbg]     (132,-64) rectangle (168,12);
  \end{scope}
  \node[font=\small\bfseries,color=accent] at (3,7)   {S1 fetch};
  \node[font=\small\bfseries,color=accent] at (77,7)  {S2 decode + execute};
  \node[font=\small\bfseries,color=accent] at (150,7) {S3 writeback};

  % --- S1 ----------------------------------------------------------------
  \node[blk,minimum width=32,minimum height=12] (fu) at (3,-6)
        {fetch\_unit\\\scriptsize PC + ibus master};
  \node[blk,minimum width=32,minimum height=12] (bp) at (3,-30)
        {branch\_predictor\\\scriptsize BTB/BHT + RAS};

  % --- pipeline registers ------------------------------------------------
  \node[draw=rulegrey,fill=white,minimum width=4,minimum height=68,inner sep=0] (ifdx) at (24,-28) {};
  \node[rotate=90,font=\scriptsize] at (24,-28) {IF/DX};
  \node[draw=rulegrey,fill=white,minimum width=4,minimum height=68,inner sep=0] (dxwb) at (130,-28) {};
  \node[rotate=90,font=\scriptsize] at (130,-28) {DX/WB};

  % --- S2 datapath, left to right ----------------------------------------
  \node[blk,minimum width=24,minimum height=12] (dec) at (42,-6)
        {decode\\\scriptsize imm\_gen};
  \node[blk,minimum width=24,minimum height=9]  (ctl) at (42,-24) {control};
  \node[blk,minimum width=24,minimum height=12] (rf)  at (76,-6)
        {regfile\\\scriptsize 32$\times$32};
  \node[blk,minimum width=24,minimum height=9]  (fwd) at (76,-24) {forward mux};

  \node[blk,minimum width=26,minimum height=10] (alu) at (112,-6)  {alu\_top / alu};
  \node[blk,minimum width=26,minimum height=10] (bu)  at (112,-22) {branch\_unit};
  \node[blk,minimum width=26,minimum height=11] (lsu) at (112,-40)
        {lsu\\\scriptsize dbus master};
  \node[blk,minimum width=26,minimum height=11] (csr) at (112,-56)
        {csr\_file\\\scriptsize + traps};

  % --- control blocks, clear of the operand rail at x = 94 ---------------
  \node[blkf,minimum width=28,minimum height=9] (hz)  at (42,-46) {hazard\_unit};
  \node[blkf,minimum width=30,minimum height=9] (exc) at (76,-46) {exception\_unit};

  % --- S3 ----------------------------------------------------------------
  \node[blk,minimum width=30,minimum height=10] (wb) at (150,-22) {writeback\_mux};

  % --- S1 into IF/DX, IF/DX into S2 --------------------------------------
  \draw[flow] (fu.east) -- (ifdx.west|-fu);
  \draw[flow] (ifdx.east|-dec) -- (dec.west);
  \draw[flow] (ifdx.east|-ctl) -- (ctl.west);

  % --- decode -> regfile -> forwarding mux -------------------------------
  \draw[flow] (dec.east) -- (rf.west);
  \draw[flow] (rf.south) -- (fwd.north);

  % --- operand distribution rail -----------------------------------------
  \draw[flow,-] (fwd.east) -- (94,-24);
  \draw[thick,draw=accent] (94,-6) -- (94,-56);
  \foreach \y in {-6,-22,-40,-56} \draw[flow] (94,\y) -- (99,\y);

  % --- S2 results into DX/WB, one arrow out to S3 ------------------------
  \draw[flow] (alu.east) -- (dxwb.west|-alu);
  \draw[flow] (bu.east)  -- (dxwb.west|-bu);
  \draw[flow] (lsu.east) -- (dxwb.west|-lsu);
  \draw[flow] (csr.east) -- (dxwb.west|-csr);
  \draw[flow] (dxwb.east|-wb) -- (wb.west);

  % --- predictor: lookup up into S1, update back along the bottom --------
  \draw[flow] (bp.north) -- (fu.south) node[lbl,midway,right=1mm]{predict};
  \draw[flow] (126,-22) -- (126,-64) -- (3,-64) -- (bp.south)
        node[lbl,pos=0.55,above]{update at commit};

  % --- redirect: leaves branch_unit downward, climbs in the gap between the
  %     operand rail and the alu column, then runs above every block --------
  \draw[flowd] (bu.north) -- (112,-14) -- (96,-14) -- (96,4) -- (3,4) -- (fu.north)
        node[lbl,pos=0.62,above]{redirect};

  % --- writeback value back into the forwarding mux ----------------------
  \draw[flow] (wb.south) -- (150,-31) -- (61,-31) -- (61,-24) -- (fwd.west);
  \node[lbl,below=0.5mm] at (78,-31) {S3$\to$S2 forwarding};
\end{tikzpicture}}

\caption{CPU block structure. Each box is a module in \file{hdl/}.}
\label{fig:cpu}
\end{figure}

\section{Top-level interface}

\begin{table}[H]
\centering
\caption{\file{cpu\_top} parameters}
\label{tab:cpuparams}
\begin{tabular}{@{}lllL{0.42\linewidth}@{}}
\toprule
\textbf{Parameter} & \textbf{Default} & \textbf{Use} & \textbf{Effect} \\
\midrule
\sig{RESET\_PC}   & \texttt{0x0000\_0000} & Reset vector & First address fetched out of reset. \\
\sig{HART\_ID}    & \texttt{0}            & \reg{mhartid} & Gives each instance its identity in a multi-core system. \\
\sig{BP\_ENTRIES} & \texttt{128}          & Predictor & BTB/BHT depth; must be a power of two. \\
\sig{RAS\_DEPTH}  & \texttt{8}            & Predictor & Return-address stack depth; \texttt{0} removes the RAS entirely. \\
\bottomrule
\end{tabular}
\end{table}

The port list is one clock, one reset, the two AXI4-Lite master ports, and the
five interrupt signals of Table~\ref{tab:cpupic}. It has not changed since the
first pipeline version; \sig{WFI}, the return-address stack and the performance
counters were all added without touching it. That was a deliberate constraint,
since the CPU is one block among several in a system being built in parallel.

\section{Datapath modules}

\begin{table}[H]
\centering
\caption{Combinational datapath modules}
\label{tab:datapath}
\begin{tabular}{@{}lL{0.66\linewidth}@{}}
\toprule
\textbf{Module} & \textbf{Function as implemented} \\
\midrule
\file{decode.v}  & Slices the fixed RV32I fields out of the instruction word: opcode, \sig{rd},
                   \sig{rs1}, \sig{rs2}, \sig{funct3}, \sig{funct7}. Pure wiring. \\
\file{imm\_gen.v} & Rebuilds the immediate for the five encoding formats from the same bits. \\
\file{control.v} & Turns the opcode into the S2 control lines, and flags any encoding it does
                   not recognise. An unknown opcode leaves every line at its inactive default;
                   a \emph{known} opcode with a bad \sig{funct} field keeps the lines its opcode
                   implies and only adds \sig{illegal}. Either way the instruction has no side
                   effect, but by two later gates rather than by the decoder: the LSU request is
                   ANDed with \sig{!illegal}, so no transaction is issued, and the trap squashes
                   the instruction on its way into DX/WB, so no register is written. Both are
                   checked in \file{tb\_traps}. \\
\file{alu\_top.v} & Selects the operands (register or immediate, PC for \sig{AUIPC}) and derives
                   the ALU operation from \sig{funct3}/\sig{funct7}. \\
\file{alu.v}     & The arithmetic and logic itself. \\
\file{branch\_unit.v} & Computes the \emph{real} direction and target of every control transfer.
                   This is the value the prediction is checked against and the value that trains
                   the predictor. Masks bit~0 of the target for \sig{JALR}, as the ISA requires. \\
\file{regfile.v} & $32\times32$, \sig{x0} hardwired to zero, two combinational read ports for S2,
                   one write port from S3, all registers reset to zero. \\
\file{writeback\_mux.v} & Selects what enters \sig{rd}: ALU result, load data, PC+4 for the
                   \sig{JAL}/\sig{JALR} link, or the CSR read value. The write enable is gated by
                   the DX/WB valid bit. \\
\bottomrule
\end{tabular}
\end{table}

\section{Instruction set coverage}

RV32I is implemented in full. Two encodings are handled by explicit decision
rather than by the base arithmetic:

\begin{itemize}[nosep]
  \item \sig{FENCE} decodes as a no-operation. The core is a single hart, in-order,
        with blocking memory operations and no store buffer, so there is nothing
        to order. Only \sig{funct3} = \texttt{000} is accepted; \sig{FENCE.I}
        belongs to Zifencei, is not implemented, and decodes as illegal.
  \item The \sig{SYSTEM} opcode decodes only exact encodings: \sig{ECALL},
        \sig{EBREAK}, \sig{MRET}, \sig{WFI} (\sig{imm12} = \texttt{0x105}) and the
        six CSR forms. \sig{funct3} = \texttt{100} is reserved and decodes as
        illegal.
\end{itemize}

Multiply and divide are out of scope per the brief and are not implemented; the
\sig{M} extension encodings decode as illegal instructions. \reg{misa} reports
\texttt{0x4000\_0100}, that is MXL = 32-bit with extension \sig{I} only, which
matches what is actually built.
